% Ad Familiares 10.8 (ad-familiares-10-08)
% Source: https://raw.githubusercontent.com/PerseusDL/canonical-latinLit/master/data/phi0474/phi056/phi0474.phi056.perseus-lat1.xml
% Fetched by scripts/fetch_latin.py

% Dateline (Perseus): Scr. in Gallia Transalpina paulo post med. Mart. a. 711 (43) .

\ciceroLetterOpener{PLANCVS IMR COS. DESIG. S. D. COS. PR. TR. PL. SENATVI POPVLO PLEBIQVE ROMANAE}

\ciceroSection{1} si cui forte videor diutius et hominum exspectationem et spem rei p. de mea voluntate tenuisse suspensam, huic prius excusandum me esse arbitror quam de insequenti officio quicquam ulli pollicendum ; non enim praeteritam culpam : videri volo redemisse, sed optimae mentis cogitata iam pridem maturo tempore enuntiare.

\ciceroSection{2} non me praeteribat in tanta sollicitudine hominum et tam perturbato statu civitatis fructuosissimam esse professionem bonae voluntatis, magnosque honores ex ea re ; compluris consecutos videbam ; sed cum in eum casum me fortuna demisisset, ut aut celeriter pollicendo magna mihi ipse ad proficiendum impedimenta opponerem aut, si in eo mihi temperavissem, maiores occasiones ad opitulandum haberem, expeditius iter communis salutis quam meae laudis esse volui. nam quis in ea fortuna quae mea est, et ab ea vita quam in me cognitam hominibus arbitror et cum ea spe quam in manibus habeo, aut sordidum quicquam pati aut perniciosum concupiscere potest ?

\ciceroSection{3} sed aliquantum nobis temporis et magni labores et multae impensae opus fuerunt ut, quae rei p. bonisque omnibus polliceremur, exitu praestaremus neque ad auxilium patriae nudi cum bona voluntate sed cum facultatibus accederemus. confirmandus erat exercitus nobis magnis saepe praemiis sollicitatus, ut a b re p. potius moderata quam ab uno infinita speraret ; confirmandae complures civitates, quae superiore anno largitionibus concessionibusque praemiorum erant obligatae, ut et illa vana putarent et eadem a melioribus auctoribus petenda existimarent; eliciendae etiam voluntates reliquorum, qui finitimis provinciis exercitibusque praefuerunt, ut potius cum pluribus societatem defendendae libertatis iniremus quam cum paucioribus funestam orbi terrarum victoriam partiremur.

\ciceroSection{4} muniendi vero nosmet ipsi fuimus aucto exercitu auxiliisque multiplicatis ut, cum praeferremus sensus aperte, tum etiam invitis quibusdam sciri quid defensuri essemus non esset periculosum. ita numquam diffitebor multa me, ut ad effectum horum consiliorum pervenirem, et simulasse invitum et dissimulasse cum dolore, quod praematura denuntiatio boni civis imparati quam periculosa esset ex casu conlegae videbam.

\ciceroSection{5} quo nomine etiam C. Furnio legato, viro forti atque strenuo, plura etiam verbo quam scriptura mandata dedimus, ut et tectius ad vos perferrentur et nos essemus tutiores, quibusque rebus et communem salutem muniri et nos armari conveniret praecepimus. ex quo intellegi potest curam rei p. summae defendendae iam pridem apud nos excubare.

\ciceroSection{6} nunc cum deum benignitate ab omni re sumus paratiores, non solum bene sperare de nobis homines sed explorate iudicare volumus. legiones habeo quinque sub signis et sua fide virtuteque rei p. coniunctissimas et nostra liberalitate nobis obsequentis, provinciam omnium civitatium consensu paratissimam et summa contentione ad officia certantem, equitatus auxiliorumque tantas copias quantas hae gentes ad defendendam suam salutem libertatemque conficere possunt ; ipse ita sum animo paratus, ut vel provinciam tueri vel ire quo res p. vocet vel tradere exercitum, auxilia provinciamque vel omnem impetum belli in me convertere non recusem, si modo meo casu aut confirmare patriae salutem aut periculum possim morari.

\ciceroSection{7} haec si iam expeditis omnibus rebus tranquilloque statu civitatis polliceor, in damno meae laudis rei p. commodo laetabor ; sin ad societatem integerrimorum et maximorum periculorum accedam, consilia mea aequis iudicibus ab obtrectatione invidorum defendenda commendo. mihi quidem ipsi fructus meritorum meorum in rei p. incolumitate satis magnus est paratus ; eos vero, qui meam auctoritatem et multo magis vestram fidem secuti nec ulla spe decipi nec ullo metu terreri potuerunt, ut commendatos vobis habeatis petendum videtur.
